\section{The equality case and the parallel splitting}
\label{sec:classification}

The preceding estimates have nonnegative integral and a nonpositive
pointwise upper bound. We now determine what their equality means.
The argument first makes the two norms constant; the norm equations
then make the Weyl tensors parallel. These are separate steps.

\begin{proposition}[Constancy of the two Weyl norms]
\label{prop:constant-norms}
On a closed, simply connected, oriented four-manifold with
$\operatorname{Ric}=3g$, nonnegative sectional curvature, and neither
Weyl block identically zero, one has $v=b=1$ everywhere.
\end{proposition}
\begin{proof}
Proposition~\ref{prop:balanced-moments} gives
$\mathbf E v^2=\mathbf E b^2$, so
Proposition~\ref{prop:coupling-estimates} supplies an integrable
function $\mathcal I$ with $\mathbf E\mathcal I\ge0$ and
\[
 \mathcal I\le-cj+2\Theta r-\frac{\Theta^2}{c},\qquad
 c\ge\frac1{10},
\]
where
\[
 r=R_v|p|^2+R_b|q|^2,\quad
 j=(|p|^2+|q|^2)\left(\frac{|p|^2}{v}+\frac{|q|^2}{b}\right),
 \quad p=dv,\quad q=db.
\]
The zero-set conventions of Proposition~\ref{prop:weak-norm-equation}
apply. The gradients are bounded, and the derivative quotients are
integrable, so $j$ is integrable as well.

For $v>0$ the identity
\[
 1+6v+v^2-4\sqrt v(1+v)=(\sqrt v-1)^4
\]
gives $R_v\le1/\sqrt v$, with equality precisely when $v=1$.
Consequently Cauchy--Schwarz gives
\[
 r\le\frac{|p|^2}{\sqrt v}+\frac{|q|^2}{\sqrt b}\le\sqrt j.
\]
A summand with zero numerator and zero denominator is read as zero.
The first inequality is strict almost everywhere on the set where
$p$ or $q$ is nonzero. To see this, recall that a Lipschitz function has
zero differential almost everywhere on each of its level sets: at a
differentiability and density point of a level set, a nonzero
differential would contradict density. Outside a null set, a nonzero
$p$ thus has $v>0$ and $v\ne1$, and the analogous statement holds for
$q$. At least one summand is then strictly smaller.

Completing the square now gives
\[
 \mathcal I\le-\frac{(\Theta-cr)^2}{c}-c(j-r^2)\le0.
\]
The right side is strictly negative almost everywhere where either
norm gradient is nonzero. Since $\mathbf E\mathcal I\ge0$, that set
has measure zero. Both weak gradients vanish; connectedness makes the
Lipschitz functions $v,b$ constant. Neither constant is zero by
hypothesis. The gaps in Proposition~\ref{prop:weyl-gap} give
$v,b\ge1$, while $v+b\le2$. Thus both constants equal $1$.
\end{proof}

The strictness in this argument uses the level sets of the Weyl norms.
It does not require the quadratic form in the Hessian identity to be
positive definite. This is why the same coupling estimates determine
the boundary geometry.

\begin{proposition}[Parallel Weyl blocks]\label{prop:parallel-blocks}
Both $T$ and $U$ are parallel, with eigenvalues $(2,-1,-1)$.
\end{proposition}
\begin{proof}
Since $v=1$ and $p=0$, the weak norm Equation \eqref{eq:3} becomes
\[
 0=6(1-\sigma)+e.
\]
Here $\sigma\le1$ and $e=|\nabla T|^2/6\ge0$. Both terms vanish,
so $\sigma=1$ and $\nabla T=0$; smoothness gives these identities
everywhere. The exchanged equation gives the same conclusion for $U$.
The trace, squared norm, and cubic trace of each block are respectively
$0,6,6$. Its characteristic polynomial is therefore
$\lambda^3-3\lambda-2=(\lambda-2)(\lambda+1)^2$, as claimed.
\end{proof}

\begin{proof}[Proof of Theorem~\ref{thm:main}]
We have reduced to the complete simply connected oriented cover.
By Proposition~\ref{prop:parallel-blocks}, the simple eigenspaces of
$T$ and $U$ are parallel real line subbundles of $\Lambda^+$ and
$\Lambda^-$. A real line bundle over a simply connected manifold is
trivial. Choose global unit sections $\omega_+$ and $\omega_-$ of
these lines. Their covariant derivatives lie in the respective lines
and are orthogonal to the unit sections, so the sections are parallel.

Viewed as skew endomorphisms, $I=\sqrt2\,\omega_+$ and
$J=\sqrt2\,\omega_-$ are orthogonal complex structures of opposite
orientations. The two rotation algebras commute, so $IJ=JI$.
The product $A=IJ$ is parallel and satisfies
\[
 A^*=A,\qquad A^2=\operatorname{Id},\qquad
 \operatorname{tr}A=0.
\]
The last equality follows from the orthogonality of self-dual and
anti-self-dual two-forms. Thus the $+1$ and $-1$ eigenspaces of $A$
are orthogonal parallel rank-two distributions.

The de Rham decomposition theorem~\cite{deRham1952} gives a global
Riemannian product of complete simply connected surfaces. The Ricci
tensor of a product restricts to the Ricci tensor of each factor, so
each surface has $\operatorname{Ric}=3g$, equivalently Gaussian
curvature $3$. A complete simply connected surface of constant
curvature $3$ is the round sphere of radius $1/\sqrt3$. This identifies
the universal Riemannian cover as the asserted product.
\end{proof}

\begin{proof}[Proof of Corollary~\ref{cor:nonnegative-classification}]
Pass to the compact, simply connected, oriented universal Riemannian
cover, as in Section~\ref{sec:introduction}. If one Weyl block vanishes
identically, the classical half-conformally-flat classification recalled
there gives the standard sphere or complex projective plane. Otherwise
Propositions~\ref{prop:constant-norms} and \ref{prop:parallel-blocks}
give parallel Weyl blocks with spectrum $(2,-1,-1)$. The parallel
splitting just proved then gives the product of two round spheres.
\end{proof}
